% Chapter 3 — read: one page becomes one structured record (kym_parse).

\gmzoomin{read}
\note{Chapter 3: reading the page the way a person would.}

\begin{frame}{Doge, step 2: read the page like a person would}
  \begin{columns}[c]
    \column{.47\textwidth}\centering
      \begin{tikzpicture}[x=1cm, y=1cm,
          blk/.style={rounded corners=1mm, fill=gm@col@read!14!gmSurface, draw=gm@col@read!60!gmSurface,
                      line width=.5pt, font=\tiny, text=gmInk, inner sep=1.5pt},
          blkd/.style={blk, fill=gm@col@read!22!gmSurface},
          blkt/.style={blk, fill=gm@col@read!30!gmSurface}]
        \draw[gmBase, line width=.8pt, rounded corners=1.5mm] (0,0) rectangle (5.6,6.3);
        \fill[gmHair] (0,5.95) rectangle (5.6,6.3);
        \foreach \x in {.18,.36,.54} \fill[gmBase] (\x,6.125) circle (.06);
        \node[blkt, minimum width=3.4cm, minimum height=.45cm, anchor=north west, font=\scriptsize\bfseries] at (.2,5.8) {Doge};
        \node[blk, minimum width=1.6cm, minimum height=1.6cm, anchor=north west] at (.2,5.2) {image};
        \node[blkd, text width=1.95cm, minimum height=1.6cm, anchor=north west, align=left] at (1.95,5.2)
          {status\\type\\year\\origin\\region};
        \node[blk, text width=1.2cm, minimum height=2.3cm, anchor=north west, align=center] at (4.2,5.8)
          {part of a series};
        \node[blkd, text width=5.05cm, minimum height=.55cm, anchor=north west] at (.2,3.45) {About};
        \node[blkd, text width=5.05cm, minimum height=.55cm, anchor=north west] at (.2,2.78) {Origin};
        \node[blkd, text width=5.05cm, minimum height=.55cm, anchor=north west] at (.2,2.11) {Spread \dots\ 21 more sections};
        \node[blk, text width=2.4cm, minimum height=.45cm, anchor=north west] at (.2,1.36) {tags};
        \node[blk, text width=2.45cm, minimum height=.45cm, anchor=north west] at (2.8,1.36) {references};
        \node[font=\tiny, text=gmMuted, anchor=south] at (2.8,.08) {knowyourmeme.com/memes/doge};
      \end{tikzpicture}
    \column{.5\textwidth}
      {\footnotesize\ttfamily
      \begin{tabular}{@{}l@{\hspace{3mm}}l@{}}
        \color{gmInk2}title    & Doge\\
        \color{gmInk2}category & meme, confirmed\\
        \color{gmInk2}year     & 2010\\
        \color{gmInk2}origin   & Tumblr\\
        \color{gmInk2}region   & Japan\\
        \color{gmInk2}types    & animal, character, exploitable,\\
                               & image macro, slang\\
        \color{gmInk2}series   & Interior Monologue Captioning\\
        \color{gmInk2}tags     & \NDogeTags\ (shiba inu, kabosu, dogecoin\dots)\\
        \color{gmInk2}sections & \NDogeSections\ --- \NDogeLinks\ links, \NDogeImages\ images\\
        \color{gmInk2}references & \NDogeRefs\\
      \end{tabular}}
  \end{columns}
  \note{
  \begin{itemize}
    \item Layman version: a person reading this page sees a title, a picture, a box of facts on
      the right, and the story underneath. The parser does exactly that, and writes it down as a
      tidy record --- one record per page, the same fields every time.
    \item Doge: a meme, confirmed by KYM's editors, from 2010, born on Tumblr, in Japan; five
      types; part of the series ``Interior Monologue Captioning''; 22 tags; 25 sections with 77
      links and 36 images; 49 references.
    \item This record is what every later step works from. Nobody reads HTML again.
  \end{itemize}}
\end{frame}

\begin{frame}{\NEntries\ records}
  \begin{columns}[c]
    \column{.38\textwidth}
      {\small\begin{tabular}{@{}r@{\hspace{2.5mm}}l@{}}
        \textbf{\NMemes}         & memes\\
        \textbf{\NCatEvents}     & events\\
        \textbf{\NCatSubcultures}& subcultures\\
        \textbf{\NCatPeople}     & people\\
        \textbf{\NCatSites}      & sites\\
        \textbf{\NCatCultures}   & cultures\\
      \end{tabular}}\par\vskip5mm
      {\small\color{gmInk2}\NReady\ have every field (\NReadyPct\%); the rest are missing something}
    \column{.58\textwidth}
      {\scriptsize\color{gmInk2}what the incomplete ones are missing}\par\vskip1mm
      % generated by scripts/figures.py — do not edit
\begin{tikzpicture}
\begin{axis}[gm hbar, scale only axis, width=.8\linewidth, height=5.86cm, xmax=23813, symbolic y coords={{region},{Spread},{Origin},{entry type},{About},{year},{tags}},
  point meta=explicit symbolic, nodes near coords={\pgfplotspointmeta}]
\addplot[fill=gm@col@read] coordinates {(18604,{region}) [18,604\enspace(77\%)] (5751,{Spread}) [5,751\enspace(24\%)] (5609,{Origin}) [5,609\enspace(23\%)] (4563,{entry type}) [4,563\enspace(19\%)] (2205,{About}) [2,205\enspace(9\%)] (770,{year}) [770\enspace(3\%)] (66,{tags}) [66\enspace(0\%)]};
\end{axis}
\end{tikzpicture}

  \end{columns}
  \note{
  \begin{itemize}
    \item 24,291 entries, all parsed by the same parser version (1.7.0).
    \item Most are memes (18,738), then events, subcultures, people, sites, cultures --- KYM's own
      categories, which become IMKG's classes (kym:Meme, kym:Event \dots).
    \item Only 18.6\% have every field. That is normal: most entries have no region, a quarter
      have no Origin or Spread section. The chart is what is missing, and how often.
    \item Important: none of them is thrown away for it --- next slide.
  \end{itemize}}
\end{frame}

\begin{frame}{A parser you can test, and a record you can trust}
  \begin{columns}[t]
    \column{.48\textwidth}
      \begin{itemize}\setlength{\itemsep}{2mm}
        \item \textbf{HTML in, record out}: no database, no scheduler inside --- tested on saved
          pages, Doge's among them
        \item \textbf{A schema with a version}: every record says which parser wrote it; a new
          parser re-reads only what it changes
        \item \textbf{Incomplete is a status, not a rejection}: the record lists what is missing
      \end{itemize}
    \column{.46\textwidth}
      \begin{tikzpicture}[baseline=(current bounding box.north)]
        \node[fill=gmTint, rounded corners=2mm, inner sep=3.5mm, text width=5.4cm, align=left] {
          \gmwhy{Pages that are not entries}\par\vskip1mm\footnotesize
          Photo galleries, videos and articles share the addresses of entries. They fail the
          schema and go to a separate ``dead letter'' list --- \NParseFailures\ of them --- so the
          entries stay clean and nothing is lost.};
      \end{tikzpicture}\par\vskip.5mm
      \begin{tikzpicture}[baseline=(current bounding box.north)]
        \node[fill=gmTint, rounded corners=2mm, inner sep=3.5mm, text width=5.4cm, align=left] {
          \gmwhy{Tried and dropped}\par\vskip1mm\footnotesize
          Grouping pages by layout automatically. An explicit parser with saved test pages was
          simpler to trust and to fix.};
      \end{tikzpicture}
  \end{columns}
  \note{
  \begin{itemize}
    \item The parser (kym\_parse.py) is a pure function: HTML in, a typed record out. No Mongo,
      no Airflow. That is why it can be tested on saved pages --- the repository keeps Doge's
      page as a test fixture.
    \item Every record carries parser\_version, schema\_version and corpus\_policy\_version. A new
      parser version re-parses; an unchanged page is not re-read (content fingerprint).
    \item Corpus policy: ``nothing is discarded for being incomplete''. Tags became gated rather
      than required (16 July). Missing fields are listed in corpus\_missing; the status is ready
      or incomplete.
    \item Non-entries go to parse\_failures (a dead-letter collection), so entries stay schema-pure.
    \item An approach I tried and removed: clustering pages by DOM structure. It found layouts,
      but an explicit parser with fixtures was easier to trust. Worth saying: dropping things is
      part of engineering.
  \end{itemize}}
\end{frame}

\gmzoomout{read}
